\documentclass[11pt]{article}
\usepackage[a4paper,margin=1in]{geometry}
\usepackage{amsmath,amssymb,amsthm,hyperref}
\newtheorem{theorem}{Theorem}
\newcommand{\C}{\mathbb C}
\newcommand{\norm}[1]{\lVert #1\rVert_R}
\title{Convergence of the Habib--Hildebrand exponential-polynomial expansion}
\author{Principia Math In Progress: MathDB 375458}
\date{22 September 2026}
\begin{document}
\maketitle

\paragraph{Task 1: full solution.}
We prove Conjecture 7.3 of Habib and Hildebrand~\cite{necessary}.
Their subsequent paper~\cite{sufficient} proves formal polynomial existence
(Theorem 2.5) and supplies the first-order system used below (Lemma 3.1),
but explicitly leaves convergence open in its final section.
These results, the quadratic example, and the standard majorant method
are credited to their respective sources; the result here is the uniform
exponential coefficient bound and its convergence consequence.

\begin{theorem}
Let $P\in\C[t]$ satisfy $P(0)=0$ and $P'(0)=P''(0)/2$.
For its formal solution
\[
 u(x)=\sum_{k\geq0}p_k(x+c)e^{-kx},\qquad p_0(t)=t,
 \qquad uu''-(u')^2=e^{-x}P(u)-1,
\]
there exists $\alpha\geq0$, depending only on $P$, such that every
coefficient $c_{k,j}=[t^j]p_k(t)$ satisfies
$|c_{k,j}|\leq e^{\alpha k}$.
For nonzero $P$, put $d=\deg P\geq2$. For every $c\in\C$ the series
converges absolutely and locally uniformly on
\[
 \mathcal D_{\alpha,c}=
 \{x\in\C:e^{\alpha-\operatorname{Re}x}
                  \max(1,|x+c|^{d-2})<1\}.
\]
Its sum is holomorphic there and solves the displayed differential equation.
For $d=2$ the factor involving $|x+c|$ is understood to be $1$.
For $P=0$ the series is simply $x+c$.
\end{theorem}

\paragraph{Polynomial recurrences and degree control.}
Write
\[
 P(t)=a(t+t^2)+t^3Q(t).
\]
Treat $t,z$ as independent formal variables and use the derivation
$\delta=\partial_t-z\partial_z$. Consider the system from
\cite[Lemma 3.1]{sufficient},
\begin{equation}\label{system}
 \delta U=1-azU+U^2V,\qquad
 \delta V=zQ(U)+azV-UV^2.
\end{equation}
Seek $U=t+\sum_{k\geq1}p_k(t)z^k$ and
$V=\sum_{k\geq1}v_k(t)z^k$.
At order $k$, the second equation determines $v_k$ from earlier
coefficients: the factor $z$ shifts the first two terms, and $V^2$ has
no terms of order less than two.
The first equation then determines $p_k$, with its only order-$k$
unknown on the right being the already determined $t^2v_k$.
Both equations have the form $(\partial_t-k)f=g$.
For a polynomial $g$ of degree at most $m$, their unique polynomial solution is
\begin{equation}\label{inverse}
 (\partial_t-k)^{-1}g
 =-\sum_{j=0}^{m}\frac{g^{(j)}}{k^{j+1}}.
\end{equation}
Thus all these polynomials exist.

Direct differentiation of \eqref{system}, using $\delta z=-z$, gives
\begin{equation}\label{identity}
 U\delta^2U-(\delta U)^2=z\{a(U+U^2)+U^3Q(U)\}-1.
\end{equation}
In particular these $p_k$ are the source's formal coefficients.
For completeness, uniqueness for the original equation also follows
directly: at order $k$, the operator on $p_k$ is
\[
 L_kp=t(p''-2kp'+k^2p)-2(p'-kp).
\]
If $p\neq0$ has leading term $bt^m$, then $L_kp$ has leading term
$k^2bt^{m+1}$, so $L_k$ is injective.

Suppose first $d\geq3$ and set $D=d-2\geq1$.
Induction in the above triangular recurrences proves
\begin{equation}\label{degree}
 \deg p_k\leq Dk+1\quad(k\geq0),\qquad
 \deg v_k\leq Dk-1\quad(k\geq1).
\end{equation}
Here zero polynomials satisfy every degree upper bound.
Indeed, since $\deg Q\leq D-1$, the coefficient of $z^k$ in
$zQ(U)$ has degree at most $D(k-1)+(D-1)=Dk-1$.
The term $azV$ has no larger degree. A product $p_\ell v_i v_j$
with $\ell+i+j=k$ and $i,j\geq1$ has degree at most $Dk-1$.
This proves the $v_k$ bound. Products $p_\ell p_m v_i$ of total
index $k$ have degree at most $Dk+1$, and $azU$ meets the same
bound. Formula \eqref{inverse} preserves degree, completing the induction.

\paragraph{A uniform inverse estimate.}
Fix $R=2(D+1)$ and, for a polynomial $f(t)=\sum_j f_jt^j$, set
\[
 \norm{f}=\sum_j|f_j|R^j.
\]
This norm is submultiplicative, $\norm{t}=R$, and for
$\deg f\leq m$ it satisfies
$\norm{f^{(j)}}\leq (m/R)^j\norm{f}$.
When $m\leq Dk+1\leq(D+1)k$, formula \eqref{inverse} therefore gives
\begin{equation}\label{bound}
 \norm{(\partial_t-k)^{-1}f}
 \leq\frac{\norm{f}}{k}
       \sum_{j=0}^{m}\left(\frac{m}{Rk}\right)^j
 \leq\frac{2}{k}\norm{f}\leq2\norm{f}.
\end{equation}
The same fixed $R$ works at every order; this is the essential estimate.

\paragraph{An analytic scalar majorant.}
Put $h=|a|$ and $Q_+(s)=\sum_j|[t^j]Q|s^j$.
Define scalar power series $A(z),B(z)$ with $A(0)=B(0)=0$ by
\begin{subequations}
\begin{align}
 A&=2\{hz(R+A)+(R+A)^2B\},\label{ma}\\
 B&=2\{zQ_+(R+A)+hzB+(R+A)B^2\}.\label{mb}
\end{align}
\end{subequations}
After moving the right sides to the left, the Jacobian with respect to
$(A,B)$ at $(z,A,B)=(0,0,0)$ is
\[
 \begin{pmatrix}1&-2R^2\\0&1\end{pmatrix}.
\]
The analytic implicit-function theorem gives a unique holomorphic
solution near $z=0$.
Its Taylor coefficients are nonnegative: at order $k$, \eqref{mb}
uses only earlier coefficients to determine $[z^k]B$; then \eqref{ma}
determines $[z^k]A$, using that coefficient of $B$ and earlier ones.
Every operation in these two recurrences has a nonnegative coefficient.

Apply the norm to the polynomial recurrences from \eqref{system},
use \eqref{bound}, and induct on $k$, determining $v_k$ first.
Submultiplicativity and the triangle inequality give
\begin{equation}\label{domination}
 \norm{v_k}\leq[z^k]B(z),\qquad
 \norm{p_k}\leq[z^k]A(z)\quad(k\geq1).
\end{equation}
For example, $zQ(U)$ is dominated by $zQ_+(R+A)$ and $U^2V$
by $(R+A)^2B$, exactly as in \eqref{ma}--\eqref{mb}.
There is no assumption that the formal polynomial series already converge.

Choose $r>0$ with $A$ holomorphic on a neighborhood of $|z|\leq r$,
and define
\[
 M=\max\bigl(1,\max_{|z|=r}|A(z)|\bigr),\qquad
 e^\alpha=\max(1,M/r).
\]
Cauchy's estimate and \eqref{domination} imply, for $k\geq1$,
\[
 |c_{k,j}|\leq\norm{p_k}\leq Mr^{-k}\leq e^{\alpha k}.
\]
The last inequality follows from $M\geq1$; it also covers $r>M$.
The assertion for $k=0$ is immediate from $p_0=t$.

\paragraph{Low degrees and convergence.}
For $P=a(t+t^2)$ the above system has $V=0$ and, as in the
explicit example of \cite[Section 4]{necessary},
\[
 p_k(t)=\frac{a^k}{k!}(t+H_k),\qquad
 H_k=\sum_{j=1}^k\frac1j,\quad H_0=0.
\]
Substitution in $(\partial_t-k)p_k=-ap_{k-1}$ verifies the formula.
Since $H_k\leq k\leq k!$ for $k\geq1$, the required coefficient
bound holds with $e^\alpha=\max(1,|a|)$; the series in fact converges
everywhere. The zero polynomial is included by taking $a=0$.

For $d\geq3$, set
$q(x)=e^{\alpha-\operatorname{Re}x}\max(1,|x+c|^D)$.
The degree and coefficient estimates imply
\[
 |p_k(x+c)e^{-kx}|
 \leq(Dk+2)\max(1,|x+c|)q(x)^k.
\]
On every compact subset of $\mathcal D_{\alpha,c}$, the continuous
function $q$ has maximum strictly below one. Thus the Weierstrass
test proves absolute local uniform convergence throughout this domain,
which is exactly the domain in \cite[Lemma 6.1]{necessary}.
Local Cauchy estimates give normal convergence of the differentiated
series, permitting differentiation of the sum. Absolute convergence also permits the
finite products in the differential equation; its formal identities
therefore become analytic identities. This proves the theorem.

\paragraph{Tasks 2--5: proposed follow-ups, not additional solved problems.}
\begin{enumerate}
 \item Determine useful explicit values of $\alpha$ from the algebraic
 majorant, and compare the resulting domains with numerical singularities.
 \item Determine the maximal convergence domain for particular polynomials
 $P$, including the monomials of degrees $3,4,6$.
 \item Resolve the separate coefficient-positivity conjecture, which the
 absolute-value majorant here does not settle.
 \item Extend the method to suitable analytic nonpolynomial functions $P$,
 where the degree estimate used here is unavailable.
\end{enumerate}

\paragraph{Verification status.}
This is an authoring-agent proof review, not an independent referee report
or a proof-assistant certificate. The companion exact-arithmetic verifier
checks the differential identity, inverse estimates, polynomial recurrences,
majorant domination, and the quadratic formula in finite examples.
The argument above, rather than those finite checks, supplies all quantifiers.

\begin{thebibliography}{9}
\bibitem{necessary} R. Habib and R. Hildebrand,
\emph{Necessary conditions for the existence of exponential-polynomial
expansions for solutions of certain differential equations},
\href{https://arxiv.org/abs/2606.21145v1}{arXiv:2606.21145v1}, 2026.
\bibitem{sufficient} R. Hildebrand and R. Habib,
\emph{Sufficient conditions for the existence of exponential-polynomial
expansions for solutions of certain differential equations},
\href{https://arxiv.org/abs/2606.24182v1}{arXiv:2606.24182v1}, 2026.
\end{thebibliography}
\end{document}
